\chapter{Introdução}

O \textit{extended Berkeley Packet Filter} (eBPF) é uma tecnologia do kernel Linux
que permite a execução de programas verificados diretamente no espaço do kernel,
sem necessidade de modificá-lo ou carregar módulos adicionais
\cite{linuxkernel2024}.
Originada do BSD Packet Filter proposto por \citeonline{mccanne1993} para captura
eficiente de pacotes em nível de usuário, a tecnologia evoluiu para suportar casos
de uso muito além da análise de rede: rastreamento de desempenho, observabilidade
de sistemas e, de forma crescente, processamento de rede de alto desempenho
\cite{gregg2019}.

No contexto de redes, o eBPF assume relevância especial em ambientes de funções
virtualizadas de rede (NFV), onde a capacidade de instrumentar e monitorar funções
de rede sem impacto perceptível sobre o desempenho é um requisito crítico
\cite{etsi2014}.
\citeonline{alves2026tcc} investigaram justamente esse cenário, comparando
ferramentas de extração de comportamento de funções virtualizadas --- trabalho que
serviu de base para o minicurso documentado neste relatório.
O repositório de código e os materiais práticos utilizados estão disponíveis em
\citeonline{plinta2026repo}.

% ----------------------------------------------------------------
\section{eBPF aplicado a redes}
% ----------------------------------------------------------------

A aplicação do eBPF em redes se dá em múltiplas camadas da pilha de comunicação,
desde o processamento de pacotes no nível do driver de rede até a instrumentação
de aplicações de espaço de usuário que implementam protocolos \cite{linuxkernel2024}.

\textbf{\textit{eXpress Data Path} (XDP).}
O XDP é o ponto de intervenção mais precoce disponível para programas eBPF na
pilha de rede do kernel Linux \cite{linuxkernel2024}.
Um programa XDP é executado diretamente no driver da interface de rede, antes que
o pacote seja alocado em um buffer \texttt{sk\_buff} e entregue à pilha de rede
convencional.
Isso permite descartar, redirecionar ou modificar pacotes com latência próxima à
do hardware, viabilizando casos de uso como mitigação de DDoS, balanceamento de
carga e encaminhamento acelerado \cite{gregg2019}.

\textbf{\textit{Traffic Control} (TC).}
O subsistema de controle de tráfego do kernel expõe \textit{hooks} eBPF tanto na
entrada quanto na saída de pacotes, após o processamento de camada 2 e com acesso
ao contexto completo do \texttt{sk\_buff} \cite{linuxkernel2024}.
Essa posição na pilha oferece mais informações sobre o pacote do que o XDP ---
incluindo metadados de cgroup e namespace --- ao custo de execução mais tardia.
O TC é amplamente utilizado em implementações de políticas de rede por identidade
em clusters Kubernetes, como as do projeto Cilium \cite{cassagnes2020}.

\textbf{Rastreamento de sockets e chamadas de sistema de rede.}
Programas eBPF conectados a \textit{tracepoints} de chamadas de sistema como
\texttt{connect}, \texttt{accept}, \texttt{sendmsg} e \texttt{recvmsg} oferecem
visibilidade sobre o comportamento de rede de cada processo individualmente,
associando fluxos de rede a PIDs, UIDs e cgroups \cite{gregg2019}.
Essa granularidade é essencial em ambientes NFV, onde múltiplas funções de rede
virtualizadas compartilham o mesmo host e precisam ser monitoradas de forma isolada
\cite{venancio2021}.

\textbf{Monitoramento passivo de latência TCP.}
O eBPF permite medir a latência de conexões TCP diretamente no kernel, sem
modificar as aplicações ou gerar tráfego adicional de sonda
\cite{research2024}.
Métricas como RTT (\textit{round-trip time}), retransmissões e janela de
congestionamento são acessíveis via \textit{tracepoints} de socket e podem ser
agregadas em histogramas por fluxo, endereço de destino ou cgroup, fornecendo
uma visão de desempenho de rede de alta fidelidade \cite{gregg2020}.

\textbf{Exportação de métricas para sistemas de monitoramento.}
Ferramentas como o \texttt{ebpf\_exporter} \cite{ebpfexporter} traduzem métricas
coletadas por programas eBPF para o formato Prometheus \cite{prometheusdocs},
permitindo sua visualização em dashboards Grafana \cite{grafanadocs} e integração
com pipelines de observabilidade já existentes.
Essa cadeia --- eBPF no kernel, Prometheus para coleta, Grafana para visualização
--- é o padrão adotado no minicurso para demonstrar o monitoramento de funções
virtualizadas de rede \cite{alves2026tcc}.

% ----------------------------------------------------------------
\section{Vantagens do eBPF para redes}
% ----------------------------------------------------------------

A adoção do eBPF em aplicações de rede é motivada por um conjunto de vantagens
que o distinguem tanto de abordagens baseadas em espaço de usuário quanto de
soluções que exigem modificação do kernel \cite{cassagnes2020}.

\textbf{Baixo overhead de monitoramento.}
Ferramentas tradicionais de monitoramento de rede --- como \texttt{sysstat}
\cite{godard2024} e contadores do sistema de arquivos \texttt{/proc}
\cite{manpages2024} --- realizam coleta periódica (\textit{polling}) com resolução
tipicamente de um segundo, perdendo eventos de curta duração e introduzindo
imprecisão temporal.
\citeonline{cassagnes2020} demonstraram que programas eBPF impõem sobrecarga
significativamente menor do que ferramentas de monitoramento convencionais,
tornando o monitoramento contínuo viável mesmo em enlaces de alta taxa de
transferência.

\textbf{Granularidade por fluxo e por processo.}
Enquanto contadores de interface de rede agregam todo o tráfego de um nó,
programas eBPF no TC ou em \textit{tracepoints} de socket associam cada pacote
ou conexão ao processo, cgroup e namespace responsável \cite{linuxkernel2024}.
Em ambientes NFV, essa granularidade permite isolar o comportamento de rede de
cada função virtualizada individualmente, sem necessidade de instrumentação da
própria função \cite{venancio2021}.

\textbf{Visibilidade sem modificação das aplicações.}
O eBPF opera como observador externo: não requer modificação do código-fonte das
funções de rede monitoradas, não impõe dependências de biblioteca e não altera o
comportamento observado \cite{gregg2019}.
Essa propriedade é especialmente valiosa em ambientes NFV com funções de terceiros,
cujo código-fonte pode não estar disponível \cite{etsi2014}.

\textbf{Processamento no kernel sem chamadas de sistema adicionais.}
Ao executar diretamente no kernel, programas eBPF eliminam as trocas de contexto
necessárias para comunicar eventos ao espaço de usuário em cada pacote ou conexão
\cite{linuxkernel2024}.
Agregações --- somas, médias, histogramas --- são calculadas em mapas BPF no
kernel e lidas pelo espaço de usuário em intervalos configuráveis, reduzindo
drasticamente o volume de dados transmitidos entre kernel e aplicação de
monitoramento \cite{gregg2020}.

\textbf{Programabilidade e extensibilidade.}
Ao contrário de contadores fixos do kernel ou de filtros estáticos de
\textit{firewall}, programas eBPF expressam lógica arbitrária --- dentro das
restrições do verificador --- e podem ser atualizados dinamicamente sem
interromper o tráfego de rede ou reiniciar serviços \cite{linuxkernel2024}.
Essa flexibilidade permite adaptar a instrumentação às necessidades específicas
de cada função virtualizada sem ciclos de desenvolvimento no kernel
\cite{alves2026tcc}.

\textbf{Portabilidade entre versões de kernel.}
O mecanismo BPF CO-RE (\textit{Compile Once -- Run Everywhere}) \cite{nakryiko2020},
implementado pela libbpf \cite{libbpf}, garante que um programa eBPF compilado
uma única vez execute corretamente em diferentes versões do kernel Linux, sem
recompilação.
Em ambientes de produção com frotas heterogêneas de servidores --- comum em
infraestruturas NFV de grande escala --- essa portabilidade reduz significativamente
o custo operacional de manutenção de ferramentas de monitoramento \cite{etsi2014}.